\section{One Agent or Many? The Architecture Debate}\label{sec:architecture}

Whether to build one agent or a team is one of the most visible design choices in agentic software engineering. We argue that, once cost is held equal, the evidence gives little reason to prefer teams as such, and strong reason to prefer arrangements in which one agent's work is checked by something it cannot influence.

\subsection{The Multi-Agent Program}

The team literature began by modeling a software company. ChatDev chained executive, programmer, reviewer and tester roles in a waterfall \citep{T1}; MetaGPT's roles passed structured documents \citep{T2}; later systems added group discussion, sprints and orchestrators with progress ledgers \citep{T3,T5,A11}. These papers reported large gains, but mostly against weak comparisons. ChatDev's quality score of 0.3953, against 0.1419 for a one-shot single-agent baseline, came at 22,949 tokens against 7,183, and no single agent received the same budget \citep{T1}. The authors of a failure taxonomy built from 1,642 annotated execution traces report that popular open-source multi-agent systems failed 41\% to 86.7\% of the time, on benchmarks that differ by system \citep{T7}.

\subsection{The Cost-Matched Challenge}

Comparisons that match budget or baseline strength largely erase the advantage. In TeamBench, a study by MIT, Google Research and Google DeepMind researchers that enforced roles with operating-system permissions, a planner--executor--verifier team was not significantly better than a solo agent: +0.5 points on average ($p = 0.20$), winning 68 of 155 tasks with one model (Gemini-3 Flash) held fixed \citep{T18}. An academic comparison of six multi-agent workflows against a single GPT-4.1 controller, over ten benchmarks and three runs, found one ahead by about three points and the other five trailing by 1.71--8.76 points \citep{T25}. A single-author meta-analysis of nine studies reports that, against budget- or sample-matched baselines, multi-agent systems average $-0.042$ in accuracy and equal or fall below the baseline in 74\% of comparable contrasts \citep{T9}. Debate fares no better: no method won more than 20\% of 36 scenarios against chain-of-thought \citep{T10}.

Study design explains most of the disagreement: studies that let the team spend more, or use a one-shot baseline, find gains, while studies that match tokens, iterations or a strong solo baseline do not \citep[see also][]{A10}. Task and model moderate the result. In Google's scaling study, multi-agent configurations ranged from +80.8\% to $-70.0\%$ relative to one agent, from decomposable financial reasoning to sequential planning \citep{H11}; in Meta's Gaia2, splitting work across sub-agents helped weaker models but not Claude 4 Sonnet's cost-normalized score \citep{A26}; and the Terminal-Bench authors, led from Stanford and the Laude Institute, conclude that ``model selection is usually more important than agent scaffold'' \citep{C18}.

\subsection{Where More Agents Do Help}

The gains that survive share a mechanism: an independent check, preferably one that executes. A separate test-writing agent raised HumanEval pass@1 from 71.3\% to 79.9\%, which its authors attribute in part to tests written alongside the code that ``can be biased by the code'' \citep{T4}, and removing only the quality-assurance role in MAGIS cut resolved issues from 13.94\% to 10.63\% \citep{T6}. The check must be grounded. In TeamBench's role-mixing pool, the language-model verifier approved 49.4\% of submissions the deterministic grader rejected ($n = 1{,}083$), and in the fixed-model ablation removing the verifier raised mean partial score by 5.5 points \citep{T18}. Handoffs also lose information: mean recall of spec-critical tokens fell from 0.79 in the planner's channel to 0.24 in the executor's tool inputs and outputs \citep{T18}, and in a comparison at near-equal cost, averaged over two Gemini-2.5 models, a designer--developer--tester waterfall implemented 27.71\% of requirements against 45.72\% for one agent and 49.48\% for a developer followed by a tester, though with Gemini-2.5-Flash alone the single agent led \citep{T12}.

In these studies, parallel builders helped only when isolation and executed tests disciplined the merge. Pairs of GPT-5 or Claude Sonnet 4.5 agents succeeded about 25\% of the time, roughly half a solo baseline, and a chat channel made no significant difference \citep{T19}. Parallel Claude Code agents passed 16\% of contradictory tasks against 62\% for a single agent given the whole task, while one serial handoff reached 54\% at 1.17 to 1.18 times the single agent's token cost \citep{T22}. Engineers in isolated git worktrees with test-gated merges, by contrast, lifted Claude Sonnet 4.5 from 57.2 to 63.3 on PaperBench at about twice the cost, while engineers sharing one workspace, kept apart only by instructions, scored 55.5, below the single agent \citep{T20}. A vendor report without controls describes the same shape: Anthropic's sixteen-agent compiler project, about \$20,000 in API costs, depended on a task verifier that is ``nearly perfect'' \citep{T16}.

\subsection{Practitioner Evidence}

\begin{sloppypar}
Open-source team methods report process, not outcomes. BMAD-METHOD, spec-kit, MetaGPT and gastown document roles, handoff files and gates, but in our reading none of the four repositories reports a measured success rate (MetaGPT's results are in its papers); gastown's README lists the unmeasured claim ``Scale comfortably to 20-30 agents'' \repofile{gastownhall/gastown}{README.md}. Their trackers record the failures the controlled studies predict: a developer agent that ticked off tasks without compiling or running tests \repo{bmad-code-org/BMAD-METHOD}{95}, specifications that create ``the illusion of work'' \repo{github/spec-kit}{75}, and a loop that declared a project complete when it was not \repo{frankbria/ralph-claude-code}{91}. MetaGPT's current default team hires no QA engineer and defaults to \texttt{run\_tests=False} \repofile{FoundationAgents/MetaGPT}{metagpt/software\_company.py}. In paperclip, an agent-company tool, an open issue reports reviewer handoffs that stall silently with no retry \repo{paperclipai/paperclip}{13532}. Swarm claims are the least supported: ruflo's own audit, while confirming that its core learning loop works, found a claimed attention speedup ``fabricated with \texttt{Math.random()} at runtime'' \repofile{ruvnet/ruflo}{docs/reviews/intelligence-system-audit-2026-05-29.md}, and one commenter writes, ``i still only have 1 claude code `agent' doing everything'' \repo{ruvnet/ruflo}{958}.
\end{sloppypar}

\subsection{Assessment}

The debate is mis-posed. The number of agents does not predict success; whether any agent is an independent, grounded check on another does. Testers that run the code, tests the builder cannot edit and test-gated merges account for most of the robust gains reviewed here, while roles that only re-read or re-plan add tokens and lose information at each handoff. Builders of such systems should justify each added agent by the check it performs, against a budget-matched single agent.
